\subsection{Methodisches Vorgehen}
\label{subsec:methodisches-vorgehen}

Die Untersuchung ist als technische Fallstudie des implementierten Templates
angelegt. Ausgangspunkt ist eine Bestandsaufnahme des Master Repository. Dabei
werden die Repository-Struktur, die enthaltenen Komponenten und die
dokumentierten Verantwortungsgrenzen erfasst. Anschließend werden die
Architekturentscheidungen und die Beziehungen zwischen Frontend, Backend,
Desktop-Integration, gemeinsam genutzten Schnittstellen, Persistenz und
Deployment betrachtet. Die Analyse beschreibt den vorhandenen Gegenstand und
trennt belegbare Projekteigenschaften von Annahmen oder noch offenen
Prüfpunkten.

Darauf aufbauend wird das Profilmodell untersucht. Im Mittelpunkt stehen die
Abgrenzung der Profile, die Aufteilung zwischen gemeinsamem Kern und
profilspezifischen Bestandteilen sowie das Modell optionaler Capabilities. Die
Betrachtung des Python-basierten Toolings ergänzt diese strukturelle Analyse
um den praktischen Entwicklungsablauf. Hierzu werden insbesondere
Projektgenerierung, Systemprüfung, Installation, Entwicklungsbetrieb, Tests,
Builds und Release-Vorbereitung als zusammenhängender Workflow betrachtet.

Für die Verifikation werden vorhandene automatisierte Tests, Build-Ergebnisse
und dokumentierte Abnahmen den jeweils untersuchten Anforderungen und
Architekturentscheidungen zugeordnet. Externe oder produktspezifische
Prüfschritte werden davon abgegrenzt. Auf dieser Evidenzbasis erfolgt die
Bewertung von Wiederverwendbarkeit, Standardisierung,
Produktivitätspotenzial, Stärken und Grenzen. Abschließend werden daraus
Kriterien für geeignete Einsatzfelder und für die Übertragung auf zukünftige
Kundenprojekte abgeleitet.

Abbildung~\ref{fig:case-study-method} fasst diese aufeinander aufbauenden
Untersuchungsschritte zusammen. Die vertikale Darstellung verdeutlicht, dass
die Bewertung und der Transfer erst auf Grundlage der zuvor erhobenen und
analysierten Projektinformationen erfolgen.

\begin{figure}[H]
  \centering
  \resizebox{\textwidth}{!}{%
  \begin{tikzpicture}[node distance=6mm]
    \node[casePrimary,minimum width=28mm] (capture) {Template erfassen};
    \node[caseBox,minimum width=31mm,right=of capture] (architecture) {Architektur analysieren};
    \node[caseBox,minimum width=31mm,right=of architecture] (profiles) {Profilmodell untersuchen};
    \node[caseBox,minimum width=28mm,right=of profiles] (tooling) {Tooling untersuchen};
    \node[caseBox,minimum width=38mm,right=of tooling] (verification) {Verifikationsdaten auswerten};
    \node[caseOptional,minimum width=35mm,right=of verification] (evaluation) {Stärken und Grenzen bewerten};
    \node[caseOptional,minimum width=35mm,right=of evaluation] (transfer) {Einsatzmöglichkeiten ableiten};

    \draw[caseFlow] (capture) -- (architecture);
    \draw[caseFlow] (architecture) -- (profiles);
    \draw[caseFlow] (profiles) -- (tooling);
    \draw[caseFlow] (tooling) -- (verification);
    \draw[caseFlow] (verification) -- (evaluation);
    \draw[caseFlow] (evaluation) -- (transfer);
  \end{tikzpicture}}
  \caption{Methodische Schritte der Fallstudie}
  \label{fig:case-study-method}
\end{figure}

